% ==============================================================================
% PREÁMBULO GENERAL Y PAQUETES DEL PROYECTO
% Ubicación: config/preamble.tex
% ==============================================================================
% Este archivo centraliza la tipografía, colores institucionales y macros base
% para que cualquier ajuste visual (h1, h2, h3) se modifique en un solo lugar.
% ==============================================================================

% ------------------------------------------------------------------------------
% 1. PAQUETES BÁSICOS Y FORMATO DE PÁGINA
% ------------------------------------------------------------------------------
\usepackage[utf8]{inputenc}            % Soporte para caracteres UTF-8
\usepackage[T1]{fontenc}               % Codificación de fuentes T1
\usepackage[spanish, es-tabla]{babel}  % Idioma español (fechas, tablas, etc.)
\usepackage[margin=2.5cm]{geometry}    % Márgenes de página equilibrados (2.5 cm)
\usepackage{parskip}                   % Espaciado moderno entre párrafos sin sangría

% ------------------------------------------------------------------------------
% 2. TIPOGRAFÍA LIMPIA Y MODERNA (Latin Modern Sans-Serif)
% ------------------------------------------------------------------------------
\usepackage{lmodern}
\renewcommand{\familydefault}{\sfdefault}

% ------------------------------------------------------------------------------
% 3. MATEMÁTICAS, SÍMBOLOS Y LISTAS
% ------------------------------------------------------------------------------
\usepackage{amsmath, amssymb, amsfonts}% Fórmulas matemáticas y símbolos
\setlength{\jot}{6pt}                  % Espaciado vertical entre líneas en align y gather
\usepackage{enumerate}                 % Listas personalizables (a), (b), etc.

% ------------------------------------------------------------------------------
% 4. COLORES INSTITUCIONALES (UCV / CIENCIAS)
% ------------------------------------------------------------------------------
\usepackage{xcolor}
\definecolor{primary}{RGB}{18, 48, 86}        % Azul marino institucional UCV
\definecolor{secondary}{RGB}{24, 110, 72}     % Verde institucional
\definecolor{darkgray}{RGB}{45, 50, 60}       % Gris oscuro para texto general
\definecolor{lightbg}{RGB}{246, 248, 250}     % Fondo tenue para tablas o cajas
\definecolor{bordercolor}{RGB}{215, 222, 230} % Borde y líneas divisorias sutiles

% ------------------------------------------------------------------------------
% 5. GRÁFICOS Y DIAGRAMAS
% ------------------------------------------------------------------------------
\usepackage{graphicx}                  % Inclusión de imágenes externas
\usepackage{tikz}                      % Diagramas vectoriales para circuitos y grafos

% ------------------------------------------------------------------------------
% 6. ENLACES Y REFERENCIAS
% ------------------------------------------------------------------------------
\usepackage{hyperref}                  % Enlaces en el documento
\hypersetup{
    colorlinks=true,
    linkcolor=primary,
    urlcolor=secondary,
    citecolor=primary
}

% ------------------------------------------------------------------------------
% 7. JERARQUÍA DE TÍTULOS EN NOTACIÓN WEB: \h1, \h2, \h3, \h4
% ------------------------------------------------------------------------------
% Permite estructurar documentos usando directamente la notación web:
%   \h1{Práctica 0: Conteo y probabilidad}  -> Título principal de práctica
%   \h2{Ejercicio 1}                       -> Título de ejercicio
%   \h3{Desarrollo y respuesta}             -> Subtítulo de ítem o sección
%   \h4{Caso particular}                   -> Sub-ítem menor
%
% Centraliza tamaños, colores y espaciados para modificarlos en un solo lugar.
% ------------------------------------------------------------------------------

% --- Configuración h1 (Título principal / Práctica o Examen) ---
\newcommand{\hOneSize}{\LARGE}
\newcommand{\hOneColor}{primary}
\newcommand{\hOneLineGap}{3.5mm}            % Separación equilibrada entre el texto y la línea divisoria
\newcommand{\hOneLineWidth}{0.65\textwidth} % Ancho de la línea divisoria
\newcommand{\hOneLineHeight}{0.8pt}         % Grosor de la línea divisoria
\newcommand{\hOneTocLevel}{section}         % Nivel en la tabla de contenidos (section o subsection)

\newcommand{\hOne}[1]{%
  \par\vspace{2.5mm}%
  \phantomsection%
  \addcontentsline{toc}{\hOneTocLevel}{#1}%
  \noindent%
  \begin{minipage}{\linewidth}
    \centering
    \setlength{\parskip}{0pt}%
    {\hOneSize\bfseries\color{\hOneColor} #1}\par
    \nointerlineskip\vspace{\hOneLineGap}%
    {\color{bordercolor}\rule{\hOneLineWidth}{\hOneLineHeight}}\par
  \end{minipage}\par
  \vspace{2.5mm}%
}

% --- Separador de semestre y portadilla para el manual de exámenes ---
\newcommand{\separadorSemestre}[1]{%
  \clearpage
  \phantomsection%
  % Entrada en el índice como encabezado de categoría visual (sin número de página)
  \addtocontents{toc}{\vspace{4mm}\noindent{\bfseries\large\color{primary} #1}\par\vspace{1.5mm}}%
  % Portadilla limpia y sobria en el cuerpo del manual de exámenes
  \thispagestyle{empty}%
  \vspace*{\fill}%
  \begin{center}%
    {\LARGE\bfseries\color{primary} #1}\par%
    \vspace{4mm}%
    {\color{bordercolor}\rule{0.5\textwidth}{1pt}}\par%
    \vspace{3mm}%
    {\normalsize\color{darkgray} Exámenes parciales del período lectivo}\par%
  \end{center}%
  \vspace*{\fill}%
  \clearpage
}
\newcommand{\tituloSemestre}[1]{\separadorSemestre{#1}}

% --- Configuración h2 (Título de ejercicio) ---
\newcommand{\hTwoSize}{\Large}
\newcommand{\hTwoColor}{secondary}
\newif\ifhTwoAddToToc
\hTwoAddToToctrue                      % Por defecto activo en prácticas para indexar ejercicios
\newcommand{\hTwoTocLevel}{subsection} % En las prácticas se indexa como subsección con su número

\newcommand{\hTwo}[1]{%
  \par\vspace{3.5mm}%
  \phantomsection%
  \ifhTwoAddToToc
    \addcontentsline{toc}{\hTwoTocLevel}{#1}%
  \fi
  \noindent{\hTwoSize\bfseries\color{\hTwoColor} #1\par}%
  \nopagebreak\vspace{1.5mm}%
  \nopagebreak
}

% --- Configuración h3 (Subtítulo de ítem, pregunta o desarrollo) ---
\newcommand{\hThreeSize}{\normalsize}
\newcommand{\hThreeColor}{primary}

\newcommand{\hThree}[1]{%
  \par\vspace{4.5mm}%
  \noindent{\hThreeSize\bfseries\color{\hThreeColor} #1\par}%
  \nopagebreak\vspace{1.5mm}%
  \nopagebreak
}

% --- Configuración h4 (Sub-ítem menor o nota puntual) ---
\newcommand{\hFourSize}{\small}
\newcommand{\hFourColor}{darkgray}

\newcommand{\hFour}[1]{%
  \par\vspace{3.5mm}%
  \noindent{\hFourSize\bfseries\color{\hFourColor} #1\par}%
  \nopagebreak\vspace{1.5mm}%
  \nopagebreak
}

% --- Despachador universal de notación web: \h1, \h2, \h3, \h4 ---
% En LaTeX, escribir \h1{...} pasa el argumento 1 como #1 y el texto como #2.
\newcommand{\h}[2]{%
  \ifcase#1\relax
    #2%
  \or
    \hOne{#2}%   % \h1{...}
  \or
    \hTwo{#2}%   % \h2{...}
  \or
    \hThree{#2}% % \h3{...}
  \or
    \hFour{#2}%  % \h4{...}
  \else
    \hThree{#2}%
  \fi
}

% --- Alias de compatibilidad y nombres alternativos ---
\let\hone\hOne
\let\htwo\hTwo
\let\hthree\hThree
\let\hfour\hFour
\let\hUno\hOne
\let\hDos\hTwo
\let\hTres\hThree
\let\tituloPractica\hOne
\let\tituloEjercicio\hTwo
\let\miniTitulo\hThree
\let\subtituloItem\hThree

% --- Conclusión pedagógica ---
% Uso: \conclusion{Explicación sencilla del resultado obtenido...}
\newcommand{\conclusionSize}{\normalsize}
\newcommand{\conclusionColor}{secondary}
\newcommand{\conclusion}[1]{%
  \par\vspace{5mm}%
  \noindent{\conclusionSize\bfseries\color{\conclusionColor} Conclusión\par}%
  \vspace{1.5mm}%
  #1\par%
}

% ------------------------------------------------------------------------------
% 8. CAJAS Y ENTORNOS DESTACADOS (NOTAS Y DEMOSTRACIONES)
% ------------------------------------------------------------------------------
% Entorno para notas teóricas destacadas con marco y fondo claro
% Uso: \begin{cajaNota} ... \end{cajaNota}
\newsavebox{\cajaNotaBox}
\newenvironment{cajaNota}{%
  \par\vspace{3.5mm}%
  \noindent%
  \begin{lrbox}{\cajaNotaBox}%
    \begin{minipage}{\dimexpr\linewidth-20pt-1.6pt\relax}%
      \small
}{%
    \end{minipage}%
  \end{lrbox}%
  \begingroup
  \setlength{\fboxrule}{0.8pt}%
  \setlength{\fboxsep}{10pt}%
  \fcolorbox{bordercolor}{lightbg}{\usebox{\cajaNotaBox}}%
  \endgroup
  \par\vspace{3.5mm}%
}
\let\cajaDemostracion\cajaNota
\let\endcajaDemostracion\endcajaNota
